\section{The Joy of Thinking}

\begin{quotex}
Descartes inherited from his father enough money for a life of study and travel. He was educated from 1604 to 1612 at the Jesuit College at La Fleche, where he received a good training not only in the humanities and mathematics, but was treated with exquisite consideration. He was allowed by Fr. Chalet, the rector and evidently a sensible man, to lie in bed in the morning --- a habit which he retained all his life, and which he regarded as above all conducive to intellectual profit and comfort.
\flright{\textsc{James R. Newman}, \emph{The World of Mathematics}}
\end{quotex}


When I was around 10 or 11, my parents would drop me off on Saturday mornings at the Boston Museum of Science while they shopped at a nearby big-box store\footnote{\url{https://en.wikipedia.org/wiki/Big-box_store}}. There, I would attend a children's lecture and then wait at the library. At that time, I was mostly interested in astronomy and mathematics. I can still remember the beautiful, blue covered, four volume anthology of \emph{The World of Mathematics}.


I was fascinated by the arcane symbolism and hidden knowledge within, much like an alchemist who tracked down a secret text. That is where I first learned of the intense pleasures of thinking. I was particular fascinated by Descartes's story from the ages of 8 to 16 at the Jesuit school. He was given the leisure of thinking time in the morning, but I was rudely awakened, force fed breakfast, and then bum rushed\footnote{\url{https://en.wiktionary.org/wiki/bum\%27s_rush}} to school.

\paragraph{Thinking as pleasure}


At my current stage of life, I have finally received Descartes's gift. As the adage goes, ``After delayed gratification comes greater reward, better profit and greater glory.'' (Sunday Adelaja) I now have insomnia, the symptom of which is that I awaken too late to go back to sleep, but too early to commence the day's activities.

So I just lie in bed and allow the thoughts to flow. I understand that a lot of people don't like to be alone with their thoughts. Perhaps they have nothing to think about so they fill the emptiness with idle entertainment and other peoples' thoughts. More unfortunate are those whose thoughts are unpleasant. Perhaps they are depressed, or in morning, or their minds are filled with anxiety, worry, feat or other negative emotions. We'll get back to this in the next section.

My suggestion is to build up a rich inner life, based on the thoughts of the greatest philosophers, sages, artists, and saints. In my case, I can ponder the deep questions of metaphysics, mathematics, and science, for which there is no end. History, literature, and poetry also provide fuel for thought if you are so inclined.

If you can remain interiorly in silence, thoughts will spontaneously arise. They will start off alone, but then in groups of related thoughts. Eventually, they become a complete article, such as this one. If the thoughts are good, they will be accompanied by a special type of pleasure.

Observe carefully how and whence the thoughts arise. We casually say, ``I think,'' without understanding what that means. Clearly, the conscious part of ourselves, the human ego, does not think in any actual sense similar to saying, ``I run,'' or ``I eat''. Rather the source of thoughts is the Self and only later does the Ego become aware of them.

Some go further and realize that the thoughts come from angels, since, as intellectual beings, that is how they communicate with us. Those thoughts are very special

\paragraph{What is thinking}


Wikipedia proposes a number of incompatible theories about the nature of Thought\footnote{\url{https://en.wikipedia.org/wiki/Thought}}. Since thought is foundational, it can be neither understood nor defined by anything simpler. Hence, the esoteric approach is to actually closely and carefully observe one's thoughts. By training, the student will begin to observe three centers within. These correspond to the traditional tripartite division of the soul into the rational, animal, and vegetative soul.

\begin{itemize}


\item the intellectual center registers, thinks, calculates, contrives, seeks out, etc.;

\item the emotive center has for its province the feelings, as well as the sensations and delicate passions;

\item the motor center directs the five senses, accumulates energy within the organism through its instinctive functions and presides through its motive functions, over the consumption of this energy.
\end{itemize}


Many, even Descartes, defined all inner conscious events as ``thought'', but, as we have just seen, that is too general. Thought, therefore, need to be reserved only to the intellectual center.

\subsubsection{Positive and Negative Parts}

The first observation one makes is that each center, including the intellectual center, has a positive and negative part. We notice some of our thoughts are positive but many others are negative. If the center is well-functioning, the two parts are in balance and are reconciled by a neutralizing force. \textbf{Boris Mouravieff} describes how the intellectual center operates. As such, there are no emotional aspects to it.

\begin{quotex}
Constructive, creative ideas originate from the positive part of the Intellectual center. Yet it is the negative part, which gauges the idea, taking so to say, the measure of it. And it is on the basis of this functional polarity that the center, as a whole, renders judgment.
\end{quotex}


\subsubsection{Sectors}


The states of being can always be further divided. So we notice next that the three centers are not so isolated from each other. Rather, they interpenetrate. Hence, for example, the intellectual center can be divided into three sectors (or six, if including the negative part). Thus, there are motor, emotional, and intellectual sectors for the intellectual center. They manifest in three ways:

\begin{itemize}


\item \textbf{Motor sector}: This gives rise to mechanical thinking. Hasty conclusions gratuitous generalizations and slogans are evidence of the motor sector of the intellectual center.

\item \textbf{Emotional sector}: When this sector is active (i.e., the ego consciousness is focused there), strong emotions arise with thinking. In the negative part, these emotions are like those described above. Another manifestation includes strong feelings of anger and threat, when confronted with thoughts that challenge one's world view. Pleasant emotions may be associated with beautiful sounding but ultimately vacuous thoughts. They give the illusion of deep thinking.

\item \textbf{Thinking sector}: This is pure thinking. Here the focus is on thinking and any pleasurable sensations are a side effect, not the primary goal.
\end{itemize}


\paragraph{Higher Intellectual Center}


Beyond the intellectual center, is the higher intellectual center. The latter is usually called ``intellectual intuition'' in contrast to the rational thinking of the ordinary intellectual center. The higher center is One, that is, it has no parts and no sectors; there is no negativity associated with it.

This is the part that understands Revelation and Symbolism, which the lower intellectual center cannot do.

\paragraph{Message to Posterity}


I will leave you with the same hope that Descartes expressed:

\begin{quotex}
I hope that posterity will judge me kindly, not only to the things which I have explained, but also to those which I have intentionally omitted so as to leave to others the pleasure of discovery.
\flright{\textsc{Rene Descartes}}
\end{quotex}



\flrightit{Posted on 2022-02-06 by Cologero}

\begin{center}* * *\end{center}

\begin{footnotesize}
\begin{sffamily}
\texttt{James on 2022-02-07 at 15:12 said:}

The stereotype of having ``shower thoughts'' is actually something I've enjoyed most days . My bathroom allows in natural light from the outside through windows near the ceiling and this creates a refreshing and almost natural visual and acoustic space to be without the noise and distraction of outside worries . It is here that I've often wondered , while watching the light flicker across the tiles , how the quest for an understanding of Quantum Gravity can be assisted by understanding the nature of spiritual gravity in the infinitesimal spaces of our own souls . Or how it can be that I might be able to adapt a particular understanding of the ``invisible forces'' of the universe into something ``visible'' similar to Dante's poetic work of rendering visible the invisible but this time through the medium of video game tropes so that it might be preserved into the next cycle . Or if the tongue-in-cheek theory that there is only one electron going forward and backward through time is a mathematical manifestation of the ``you were everyone who ever lived'' premise of Andy Weir's ``The Egg'' .

Josef Pieper was right : leisure truly is the basis of culture .

\end{sffamily}
\end{footnotesize}

